\section{General real planes: a quantitative interpretation}
\label{sec:realinputs}
iPRA is a first-order theory of natural numbers. It does not literally
quantify over every subset of the real numbers, every real Cauchy
sequence, or every real plane. Equation (\ref{eq:ipraendpoint}) is its
precise exact-arithmetic theorem. Here we explain why this proves the
stated geometric result, and how to read it constructively for supplied
real-number names. The distinction is substantive, not a change of the
optimization problem.

\subsection{Uniform geometric margins}
Every unscaled coordinate model in this report satisfies
\[
 \tfrac12\mathbb B\subseteq K\subseteq2\mathbb B,
                                                               \tag{R1}
\]
where $\mathbb B$ is the Euclidean unit ball. The outer inclusion follows
by checking squared vertex norms against $4$. For the inner inclusion,
orient each facet inequality as $v\cdot x\le d$ with $d>0$ and check
$4d^2\ge v\cdot v$. These are finite checks in $K=\QQ(\varphi)$ and
are also performed by the witness program. Convexity gives, for
$0<\epsilon<1$,
\begin{equation}
 (1-\epsilon)K+\tfrac\epsilon2\mathbb B\subseteq K.
                                                               \label{eq:margin}
\end{equation}
Every planar section is contained in a disk of radius $2$, hence has
area at most $16$ by an enclosing square. In particular each displayed
raw maximum $M$ is at most $16$.

\subsection{Approximate a plane after contracting its section}
Let $S=K\cap H(u,t)$, with $\|u\|=1$. If $S$ is nonempty then
$|t|\le2$. Fix a rational $0<\epsilon\le1/2$, and let $s=1-\epsilon$.
Choose a nonzero dyadic vector $n$ so that, on putting $w=n/\|n\|$,
\[
 \|w-u\|\le\epsilon/8.
\]
This is an explicit finite approximation: component errors at most
$\epsilon/64$ suffice, since $\|n-u\|\le\epsilon/32$ and
$\|n/\|n\|-u\|\le2\|n-u\|$. In particular $\|n\|\ge1/2$.
Choose a dyadic $d$ with
\[
 |d-\|n\|st|\le\epsilon/8,
 \qquad t'=d/\|n\|,
\]
so $|t'-st|\le\epsilon/4$.

There is a rotation $R$ taking $u$ to $w$ with
$\|R-\mathrm{id}\|_{\mathrm{op}}=\|u-w\|$.
For a finite algebraic verification, when $u\cdot w>-1$ put
$v=u\times w$ and use
\[
 R=\mathrm{id}+[v]_\times+\frac{[v]_\times^2}{1+u\cdot w},
 \qquad [v]_\times x=v\times x.
\]
Matrix multiplication verifies orthogonality and $Ru=w$; restriction
to the plane spanned by $u,w$ verifies the displayed operator-norm
identity. In our application the denominator is bounded away from
zero. If $u=w$, the identity rotation suffices and the same formula
already gives it.

The congruent copy
\[
 S'=R(sS)+(t'-st)w
\]
lies in the dyadic plane $H(n,d)$ and has area $s^2\Area(S)$.
For every $x\in S$,
\[
 \|R(sx)+(t'-st)w-sx\|
 \le2\|R-\mathrm{id}\|_{\mathrm{op}}+|t'-st|
 \le\epsilon/2.
\]
By (\ref{eq:margin}), $S'\subseteq K$. Therefore the proved
arithmetic-plane inequality gives
\[
 s^2\Area(S)\le\Area(K\cap H(n,d))\le M,
 \qquad \Area(S)\le\frac{M}{(1-\epsilon)^2}.        \tag{R2}
\]
Since $\epsilon\le1/2$,
\[
 \frac1{(1-\epsilon)^2}-1\le8\epsilon.
\]
Taking $\epsilon=2^{-(k+8)}$ in (R2) proves the explicit bound
\[
 \Area(S)\le M+2^{-k}
 \qquad\hbox{for every }k\in\mathbb N.             \tag{R3}
\]
The constants above actually leave a factor of two in reserve.
This argument includes supporting faces. It does not assert that
section area is continuous across a supporting plane, which would be
false because the nearby exterior sections are empty.

\subsection{Classical and constructive readings}
In ordinary real geometry, (R3) for every $k$ is exactly the sharp
inequality $\Area(S)\le M$. Together with the explicit attaining
algebraic planes, it proves the optimization result over all real
affine planes, not merely planes with algebraic coefficients.

For a constructive reading, supply a unit normal and offset by Cauchy
names with a specified error modulus, or supply a nonzero normal
with a positive lower bound on its norm and then normalize it. For
each requested accuracy, the construction queries finitely many
approximations and produces dyadic $n,d$ and the quantitative
implication (R2). Square roots of positive rationals are approximated
by bounded integer bisection; the required reciprocal has a supplied
positive denominator bound. These are primitive recursive
operations. The resulting family of finite precision statements
has an iPRA proof for every fixed primitive recursive family of
input names whose approximation contracts are available in iPRA.
Here the area upper bound is interpreted as a predicate on finite
inscribed polygon data, or conditionally on a supplied area name with
its geometric approximation contract. An area name for the full
section is not computed from arbitrary plane names: such a uniform
computation would fail at supporting planes. The construction computes
the dyadic comparison plane. Its upper bound applies to every finite
inscribed polygon, and hence to a supplied section-area name satisfying
the stated contract. The inequality between represented real areas is
read through these finite precision statements.

This is a schema for specified names and their contracts. It is not
an internally quantified claim that every real is coded by a natural
number, nor a total universal evaluator for arbitrary purported
primitive recursive program codes. The exact theorem, its
constructive finite precision consequences, and its ordinary-real
interpretation are thus kept distinct and connected by an explicit
error calculation.
